% TOP_PROBLEM: {"id": 3375, "title": "Are there infinite number of Mersenne Primes?", "category_id": 1, "category": {"id": 1, "name": "number_theory", "display_name": "Number Theory", "description": "Properties of integers, prime numbers, Diophantine equations.", "slug": "number-theory", "order_index": 1, "created_at": "2026-07-31T15:26:25.670Z"}, "status": "open", "problem_number": "OPG-59977", "set_id": 12}
% FIRST_POSTED: 2026-10-01
% ENTRY_KIND: statement_only
% UnsolvedMath ID 3375; source catalogue code OPG-59977
% !TeX program = lualatex
% Definition and sources only; this file is not an LLM solution attempt.
\documentclass[11pt,a4paper]{article}
\usepackage[margin=25mm]{geometry}
\usepackage{fontspec}
\setmainfont{lmroman10-regular.otf}[BoldFont=lmroman10-bold.otf,
 ItalicFont=lmroman10-italic.otf,BoldItalicFont=lmroman10-bolditalic.otf]
\usepackage{amsmath,amssymb,mathtools}
\usepackage{unicode-math}
\setmathfont{latinmodern-math.otf}
\AtBeginDocument{\let\setminus\smallsetminus}
\usepackage{xurl,hyperref}
\hypersetup{colorlinks=true,urlcolor=blue}
\setcounter{secnumdepth}{0}
\setlength{\parindent}{0pt}
\setlength{\parskip}{5pt}
\setlength{\emergencystretch}{3em}
\begin{document}
\section*{Are there infinite number of Mersenne Primes?}\label{3375}
{\small UnsolvedMath ID 3375 · OPG-59977 · Number Theory · OpenGarden}
\subsection{Definitions and mathematical statement}
A Mersenne number is an integer of the form $2^n-1$, where $n\ge2$.
When it is prime it is called a Mersenne prime. If $n=ab$ with
$a,b>1$, then $2^a-1$ is a proper divisor of $2^n-1$, so a prime
Mersenne number necessarily has a prime exponent. Consequently the
infinitude question can be stated exactly as
\[
 \forall X>0\ \exists\text{ prime }p>X\quad 2^p-1\text{ is prime}.
\]
Equivalently the counting function
$\#\{p\le x:p\text{ prime and }2^p-1\text{ prime}\}$ tends to
infinity as $x\to\infty$. No upper bound on the gap between successful
exponents is prescribed.

The primality of $p$ is necessary but not sufficient for the primality
of $2^p-1$. The target concerns the latter integers themselves, not
the number of prime divisors appearing among them. Either a proof of
infinitely many successful exponents or a proof that all sufficiently
large exponents fail would resolve the question. A finite collection
of certified examples establishes only the existence of that many
Mersenne primes, regardless of how large the examples are.
\subsection{Short English statement}
Are there infinitely many prime exponents p for which two to the pth power minus one is also prime?
\subsection{Sources}
\begin{itemize}
\item[S1] ULAM.AI and the UnsolvedMath Contributors, supplied problems.json, record 3375 (OPG-59977), OpenGarden collection. Catalogue identity and source transcription; CC BY 4.0.
\url{https://huggingface.co/datasets/ulamai/UnsolvedMath}.
\item[S2] Open Problem Garden, Are there infinite number of Mersenne Primes?. Original problem and discussion; the mathematical formulation below is expanded from the supplied source record.
\url{http://www.openproblemgarden.org/op/are_there_infinite_number_of_mersenne_primes}.
\end{itemize}
\end{document}
